%==============================================================
%  lecons/algebre/produit_scalaire.tex
%  Produit scalaire — espaces euclidiens
%==============================================================

\lecontitle{Produit scalaire}{Algèbre linéaire — Espaces euclidiens}

%=============================================================
%  § 1 — DÉFINITION ET EXEMPLES
%=============================================================
\secnum{navyblue}{1}{Définition et exemples}

\begin{tcolorbox}[thmbox]
  \textbf{Définition.}\enspace%
  \index{produit scalaire}\index{espace euclidien}
  Soit $E$ un $\mathbb{R}$-espace vectoriel. Un \emph{produit scalaire} sur $E$
  est une application $\langle\cdot,\cdot\rangle : E \times E \to \mathbb{R}$
  qui est :
  \begin{enumerate}
    \item \textbf{Bilinéaire :} linéaire en chaque variable ;
    \item \textbf{Symétrique :} $\langle x,y\rangle = \langle y,x\rangle$ ;
    \item \textbf{Définie positive :} $\langle x,x\rangle \geq 0$ et
          $\langle x,x\rangle = 0 \Rightarrow x = 0$.
  \end{enumerate}
  Un $\mathbb{R}$-espace vectoriel de dimension finie muni d'un produit
  scalaire est appelé \emph{espace euclidien}.
\end{tcolorbox}

\begin{mdframed}[style=exbar]
  \textbf{1. Produit scalaire canonique sur $\mathbb{R}^n$.}\enspace%
  \index{produit scalaire!canonique}
  $\langle x,y\rangle = \displaystyle\sum_{i=1}^n x_i y_i = x^T y$.
  C'est le produit scalaire usuel de la géométrie.

  \smallskip\noindent
  \textbf{2. Produit scalaire $L^2$ sur $\mathcal{C}([a,b])$.}\enspace%
  \index{produit scalaire!intégral}
  $\langle f,g\rangle = \displaystyle\int_a^b f(t)\,g(t)\,\mathrm{d}t$.
  La définition positive résulte du fait qu'une fonction continue positive
  d'intégrale nulle est identiquement nulle.

  \smallskip\noindent
  \textbf{3. Produit scalaire sur $\mathcal{M}_n(\mathbb{R})$.}\enspace%
  $\langle A,B\rangle = \mathrm{tr}(A^T B)$.
  La bilinéarité et la symétrie sont immédiates ; la définition positive
  vient de $\mathrm{tr}(A^T A) = \sum_{i,j} a_{ij}^2$.
\end{mdframed}

\decsep{}

%=============================================================
%  § 2 — INÉGALITÉ DE CAUCHY-SCHWARZ
%=============================================================
\secnum{navyblue}{2}{Inégalité de Cauchy-Schwarz}

\begin{tcolorbox}[thmbox]
  \textbf{Théorème (Cauchy-Schwarz).}\enspace%
  \index{inégalité de Cauchy-Schwarz}
  Pour tous $x, y \in E$ :
  \[
    \langle x, y\rangle^2 \leq \langle x,x\rangle\,\langle y,y\rangle.
  \]
  L'égalité a lieu si et seulement si $x$ et $y$ sont colinéaires.
\end{tcolorbox}

\begin{mdframed}[style=proofbar]
  Si $x = 0$, les deux membres sont nuls. Supposons $x \neq 0$.
  Pour tout $\lambda \in \mathbb{R}$, la définition positive donne :
  \[
    0 \leq \langle x - \lambda y,\, x - \lambda y\rangle
    = \langle x,x\rangle - 2\lambda\langle x,y\rangle + \lambda^2\langle y,y\rangle.
  \]
  C'est un trinôme en $\lambda$ positif ou nul, donc son discriminant est $\leq 0$ :
  \[
    4\langle x,y\rangle^2 - 4\langle x,x\rangle\langle y,y\rangle \leq 0.
  \]
  L'égalité a lieu si et seulement si le trinôme a une racine, i.e.\ si
  $\langle x - \lambda_0 y, x - \lambda_0 y\rangle = 0$ pour un $\lambda_0$,
  soit $x = \lambda_0 y$. \hfill$\blacksquare$
\end{mdframed}

\rubriq{Conséquences directes}
\begin{mdframed}[style=exbar]
  \textbf{Sur $\mathbb{R}^n$ :}\enspace%
  $\left(\sum x_i y_i\right)^2 \leq \left(\sum x_i^2\right)\left(\sum y_i^2\right)$.

  \smallskip\noindent
  \textbf{Sur $\mathcal{C}([a,b])$ :}\enspace%
  $\left(\int_a^b fg\right)^2 \leq \left(\int_a^b f^2\right)\left(\int_a^b g^2\right)$.
\end{mdframed}

\decsep{}

%=============================================================
%  § 3 — NORME EUCLIDIENNE
%=============================================================
\secnum{navyblue}{3}{Norme euclidienne et identités}

\begin{tcolorbox}[thmbox]
  \textbf{Théorème.}\enspace%
  L'application $\|\cdot\| : E \to \mathbb{R}_+$,\enspace%
  $\|x\| = \sqrt{\langle x,x\rangle}$, est une \emph{norme}\index{norme euclidienne}
  sur $E$, appelée \emph{norme euclidienne}. En particulier :
  \[
    \|x + y\| \leq \|x\| + \|y\| \quad\textit{(inégalité triangulaire)}.
  \]

  \medskip
  \textbf{Identités remarquables.}
  \begin{enumerate}
    \item \textit{Parallélogramme :}
    $\|x+y\|^2 + \|x-y\|^2 = 2(\|x\|^2 + \|y\|^2)$.
    \item \textit{Polarisation :}
    $\langle x,y\rangle = \dfrac{1}{4}\bigl(\|x+y\|^2 - \|x-y\|^2\bigr)$.
  \end{enumerate}
\end{tcolorbox}

\begin{mdframed}[style=proofbar]
  \textit{Inégalité triangulaire.}
  \begin{align*}
    \|x+y\|^2 &= \langle x+y,x+y\rangle = \|x\|^2 + 2\langle x,y\rangle + \|y\|^2 \\
              &\leq \|x\|^2 + 2\|x\|\,\|y\| + \|y\|^2 = {(\|x\|+\|y\|)}^2,
  \end{align*}
  où on a utilisé Cauchy-Schwarz : $|\langle x,y\rangle| \leq \|x\|\,\|y\|$.

  \textit{Identité du parallélogramme.}
  $\|x+y\|^2 + \|x-y\|^2 = 2\|x\|^2 + 2\|y\|^2$ par développement direct.

  \textit{Polarisation.} Découle immédiatement de l'identité du parallélogramme.
  \hfill$\blacksquare$
\end{mdframed}

\rubriq{Réciproque — caractérisation des normes euclidiennes}
Une norme vérifie l'identité du parallélogramme si et seulement si elle provient
d'un produit scalaire (via la formule de polarisation). La norme $\|\cdot\|_1$
sur $\mathbb{R}^2$ ne vérifie pas cette identité : elle n'est pas euclidienne.

\decsep{}

%=============================================================
%  § 4 — ORTHOGONALITÉ
%=============================================================
\secnum{navyblue}{4}{Orthogonalité et supplémentaire orthogonal}

\begin{tcolorbox}[thmbox]
  \textbf{Définition.}\enspace%
  \index{orthogonalité}\index{complement orthogonal@complément orthogonal}
  On dit que $x$ et $y$ sont \emph{orthogonaux}, noté $x \perp y$, si
  $\langle x,y\rangle = 0$.
  Le \emph{complément orthogonal} d'une partie $A \subset E$ est
  \[
    A^\perp = \bigl\{x \in E \mid \forall\,a \in A,\;\langle x,a\rangle = 0\bigr\}.
  \]
  C'est toujours un sous-espace vectoriel de $E$.

  \medskip
  \textbf{Théorème.}\enspace%
  Si $F$ est un sous-espace de l'espace euclidien $E$, alors
  \[
    E = F \oplus F^\perp, \qquad \dim F^\perp = \dim E - \dim F.
  \]
\end{tcolorbox}

\begin{mdframed}[style=proofbar]
  Soit $(e_1,\ldots,e_p)$ une base orthonormale de $F$ (obtenue par Gram-Schmidt,
  voir § 5). Pour $x \in E$, posons $p_F(x) = \sum_{i=1}^p \langle x, e_i\rangle\, e_i$.
  Alors $p_F(x) \in F$ et $x - p_F(x) \in F^\perp$ (vérification : pour tout $j$,
  $\langle x - p_F(x), e_j\rangle = \langle x,e_j\rangle - \langle x,e_j\rangle = 0$).
  Donc $x = p_F(x) + (x-p_F(x)) \in F + F^\perp$.
  Si $y \in F \cap F^\perp$, alors $\langle y,y\rangle = 0$ d'où $y = 0$.
  La somme est donc directe. \hfill$\blacksquare$
\end{mdframed}

\rubriq{Projection orthogonale}

\begin{tcolorbox}[thmbox]
  \textbf{Définition/Théorème.}\enspace%
  \index{projection orthogonale}
  L'endomorphisme $p_F : E \to E$ défini par la décomposition
  $x = p_F(x) + (x - p_F(x)) \in F \oplus F^\perp$
  est appelé \emph{projection orthogonale} sur $F$.
  C'est le point de $F$ le plus proche de $x$ :
  \[
    \|x - p_F(x)\| = \min_{y \in F} \|x - y\|.
  \]
\end{tcolorbox}

\begin{mdframed}[style=proofbar]
  Pour $y \in F$ : $\|x-y\|^2 = \|x - p_F(x) + p_F(x) - y\|^2$.
  Or $x - p_F(x) \in F^\perp$ et $p_F(x) - y \in F$, donc ils sont orthogonaux :
  \[
    \|x-y\|^2 = \|x-p_F(x)\|^2 + \|p_F(x)-y\|^2 \geq \|x-p_F(x)\|^2,
  \]
  avec égalité si et seulement si $y = p_F(x)$. \hfill$\blacksquare$
\end{mdframed}

\decsep{}

%=============================================================
%  § 5 — GRAM-SCHMIDT ET BASES ORTHONORMÉES
%=============================================================
\secnum{navyblue}{5}{Orthonormalisation de Gram-Schmidt}

\begin{tcolorbox}[thmbox]
  \textbf{Théorème (Gram-Schmidt).}\enspace%
  \index{Gram-Schmidt}\index{base orthonormale}
  Soit $(x_1, \ldots, x_n)$ une base d'un espace euclidien $E$.
  On construit par récurrence une base orthonormale $(e_1,\ldots,e_n)$ telle que
  $\mathrm{Vect}(e_1,\ldots,e_k) = \mathrm{Vect}(x_1,\ldots,x_k)$ pour tout $k$ :
  \[
    e_1 = \frac{x_1}{\|x_1\|}, \qquad
    \tilde{e}_k = x_k - \sum_{i=1}^{k-1}\langle x_k, e_i\rangle\,e_i, \qquad
    e_k = \frac{\tilde{e}_k}{\|\tilde{e}_k\|}.
  \]
  Tout espace euclidien admet une base orthonormale.
\end{tcolorbox}

\begin{mdframed}[style=proofbar]
  Par récurrence. Au rang $k$, $\tilde{e}_k$ est $x_k$ privé de sa projection
  sur $\mathrm{Vect}(e_1,\ldots,e_{k-1})$. Donc $\tilde{e}_k \perp e_i$ pour
  $i < k$ (par construction), et $\tilde{e}_k \neq 0$ car $x_k \notin
  \mathrm{Vect}(x_1,\ldots,x_{k-1})$ (famille libre). En normalisant, $e_k$
  complète la famille orthonormale. \hfill$\blacksquare$
\end{mdframed}

\rubriq{Identité de Parseval}

\begin{tcolorbox}[thmbox]
  \textbf{Théorème (Parseval).}\enspace%
  \index{identité de Parseval}
  Soit $(e_1,\ldots,e_n)$ une base orthonormale de $E$. Pour tous $x, y \in E$ :
  \[
    \langle x,y\rangle = \sum_{i=1}^n \langle x,e_i\rangle\,\langle y,e_i\rangle,
    \qquad
    \|x\|^2 = \sum_{i=1}^n \langle x, e_i\rangle^2.
  \]
\end{tcolorbox}

\begin{mdframed}[style=proofbar]
  On écrit $x = \sum_i \langle x,e_i\rangle\,e_i$ et
  $y = \sum_j \langle y,e_j\rangle\,e_j$ dans la base orthonormale.
  La bilinéarité et $\langle e_i,e_j\rangle = \delta_{ij}$ donnent le résultat.
  \hfill$\blacksquare$
\end{mdframed}

\decsep{}

%=============================================================
%  § 6 — PIÈGES ET EXERCICES
%=============================================================
\secnum{exgreen}{6}{Pièges et exercices}

\noindent{\bfseries\color{counterred} Pièges.}

\begin{mdframed}[style=cexbar]
  \textbf{Norme sans produit scalaire.}\enspace%
  Toute norme euclidienne vérifie l'identité du parallélogramme, mais
  $\|\cdot\|_1$ et $\|\cdot\|_\infty$ sur $\mathbb{R}^2$ ne la vérifient pas :
  elles ne proviennent d'aucun produit scalaire.

  \smallskip\noindent
  \textbf{$F^\perp$ dépend du produit scalaire.}\enspace%
  Si on change de produit scalaire sur $E$, le complément orthogonal d'un
  sous-espace $F$ change. L'orthogonalité n'est pas une notion purement vectorielle.

  \smallskip\noindent
  \textbf{$E = F \oplus F^\perp$ en dimension infinie.}\enspace%
  Ce résultat est faux en général dans un espace de dimension infinie non complet.
  Il tient dans les espaces de Hilbert (complets) mais pas dans tous les
  espaces préhilbertiens.
\end{mdframed}

\vspace{6pt}
\exolabel

\begin{mdframed}[style=exobar]
  \begin{enumerate}
    \item Montrer que $\langle f,g\rangle = \int_0^1 f(t)\,g(t)\,\mathrm{d}t$
          est un produit scalaire sur $\mathcal{C}([0,1],\mathbb{R})$.
          Appliquer Cauchy-Schwarz pour montrer
          $\left(\int_0^1 f\right)^2 \leq \int_0^1 f^2$.

    \item Orthonormaliser la famille $(1, X, X^2)$ dans $\mathbb{R}[X]_{\leq 2}$
          muni de $\langle P,Q\rangle = \int_{-1}^1 PQ$. Reconnaître les
          polynômes obtenus (polynômes de Legendre\index{polynômes de Legendre}).

    \item Soit $p$ la projection orthogonale sur $F = \mathrm{Vect}(e_1, e_2)$
          dans $\mathbb{R}^4$ muni du produit scalaire canonique, avec
          $e_1 = \frac{1}{2}(1,1,1,1)$, $e_2 = \frac{1}{2}(1,1,-1,-1)$.
          Calculer $p(1,0,0,0)$ et $\mathrm{dist}((1,0,0,0), F)$.

    \item Montrer que si $A$ est symétrique définie positive (i.e.\
          $A^T = A$ et $x^TAx > 0$ pour $x\neq 0$), alors
          $\langle x,y\rangle_A = x^T A y$ est un produit scalaire sur $\mathbb{R}^n$.

    \item (Inégalité de Bessel) Soit $(e_1,\ldots,e_p)$ une famille
          orthonormale dans $E$ (pas nécessairement une base).
          Montrer que $\sum_{i=1}^p \langle x,e_i\rangle^2 \leq \|x\|^2$,
          avec égalité si et seulement si $x \in \mathrm{Vect}(e_1,\ldots,e_p)$.
  \end{enumerate}
\end{mdframed}

\leconfooter{A.-L.\ Cauchy (1821) — H.\ A.\ Schwarz (1885) — J.\ P.\ Gram, E.\ Schmidt (1907)}
